% Generated draft: requires mathematical and source review.
\documentclass[11pt]{article}
\usepackage[margin=25mm]{geometry}
\usepackage{fontspec}
\setmainfont{Latin Modern Roman}
\usepackage{amsmath,amssymb,mathtools,unicode-math}
\setmathfont{Latin Modern Math}
\usepackage{enumitem,needspace,xurl,hyperref}
\setlength{\parindent}{0pt}
\setlength{\parskip}{5pt}
\setlength{\emergencystretch}{3em}
\newcommand{\N}{\mathbb{N}}
\newcommand{\Z}{\mathbb{Z}}
\newcommand{\Q}{\mathbb{Q}}
\newcommand{\R}{\mathbb{R}}
\newcommand{\C}{\mathbb{C}}
\newcommand{\F}{\mathbb{F}}
\newcommand{\A}{\mathbb{A}}
\newcommand{\PP}{\mathbb{P}}
\newcommand{\E}{\mathbb{E}}
\newcommand{\eps}{\varepsilon}
\newcommand{\dd}{\,\mathrm{d}}
\begin{document}
\section*{UnsolvedMath ID: 30000549}
\textbf{Source problem code:} catalogue entry
\subsection{Definitions and mathematical statement}
Let $\Sigma_g$ and $\Sigma_h$ denote closed connected orientable surfaces of genera $g,h$. The intended nontrivial surface-group setting normally has $g,h\ge2$, but the extraction does not state this range explicitly; it must be confirmed because allowing genus zero changes the existence question.

The source asks whether a word-hyperbolic group $\Gamma$ can fit into a short exact sequence
\[
 1\longrightarrow\pi_1(\Sigma_g)
   \xrightarrow{\iota}\Gamma
   \xrightarrow{q}\pi_1(\Sigma_h)
   \longrightarrow1.
\]
Exactness means $\iota$ is injective, $q$ is surjective, and $\ker q=\operatorname{im}\iota$. No splitting of the extension is assumed. Word hyperbolicity means that a Cayley graph for a finite generating set has uniformly thin geodesic triangles; no realization as a negatively curved manifold is required.

There are two distinct questions:
\begin{enumerate}
\item Does such an extension exist in the intended closed higher-genus setting? The bare existence wording does not specify whether $g,h$ are fixed inputs or whether some pair is sought. At minimum it asks for existence of an example; existence for every prescribed pair is a stronger assertion not silently added.
\item If such groups exist, for each fixed admissible pair $(g,h)$ are there only finitely many such groups? With the usual group-isomorphism interpretation, this asks whether
\[
 \#\{[\Gamma]_{\cong}:\Gamma\text{ is word-hyperbolic and admits the displayed extension}\}<\infty.
\]
This counts abstract groups, not marked injections, quotient maps, or extension equivalence classes; the exact counting convention requires source confirmation.
\end{enumerate}

The supplied assessment distinguishes known free-quotient or special convex-cocompact constructions from the requested closed-surface quotient. Such examples are not substituted for the stated quotient group. No uniform bound as the genera vary, no semidirect-product hypothesis, and no specific monodromy representation is added.
\subsection{Short English statement}
The questions ask whether a word-hyperbolic group can be an extension of one closed surface group by another, and whether only finitely many such group types occur for fixed genera. The genus range, existence quantifier, and equivalence used for counting must be checked.
\subsection{Sources}
\begin{itemize}
\item Local UnsolvedMath registry, permanent ID 30000549, code catalogue entry, ``Hyperbolic Surface-by-Surface Group Extensions.'' Proposed year 2006; proposer unspecified.
\item \emph{Geometric Group Theory, Hyperbolic Dynamics and Symplectic Geometry} (2007), Oberwolfach Reports, pp.\ 1991--2058, DOI 10.4171/owr/2006/33; \url{https://doi.org/10.4171/owr/2006/33}. The displayed report year 2007 differs from the DOI and proposed year 2006.
\item Supplied ``Farb--Mosher, convex-cocompact subgroups'': \url{https://arxiv.org/abs/math/0106190}.
\item Statement status corrected\_verified; the supplied note says editorial and LaTeX cleanup preserved meaning and added no source-dependent content.
\item Literature triage checked 2026-08-21: open. It describes both existence and the proposed fixed-genus finiteness as open and says known constructions with free or special quotients do not settle the closed-surface quotient. Exact hypotheses and current status require independent verification. Created 2026-08-25; updated 2026-09-16.
\item Registry curation by Ulam AI / UnsolvedMath Contributors, CC BY 4.0; underlying report and literature provenance remain as supplied. Unreviewed statement-only draft.
\end{itemize}
\end{document}
